\chapter{Introduction}
\label{ch:intro}

\section{Background}
\label{sec:background}

Vision-based robot learning benefits from demonstrations that cover different
objects, robot embodiments, camera viewpoints, and scene conditions. Direct
teleoperation provides robot actions paired with observations, but acquiring
additional recordings requires operator time and access to the relevant
hardware. The cost depends on the collection interface and task; it need not
grow in a fixed proportion to the number of robots. Demonstration reuse and
synthetic data generation offer ways to reduce repeated collection
\cite{mimicgen, real2render2real}.

Human demonstrations provide another source of manipulation motion. A person
can perform a task while cameras record hand--object interaction, without
controlling a robot during capture. Such recordings do not directly supply
robot joint states or executable actions. Their use requires estimating motion,
identifying task-relevant interactions, and adapting the resulting trajectory
to the target robot \cite{dexmv, okami}.

3D Gaussian Splatting (3DGS) supports novel-view synthesis from calibrated
images and provides an editable representation of scene appearance
\cite{gaussiansplatting}. Prior systems already use this representation for
robot-data augmentation and human-to-robot data generation
\cite{robosplat, real2render2real, robopaint}.
This thesis investigates a ViKi workflow that combines markerless human-motion
reconstruction from multiview RGB-D recordings with embodiment-specific
retargeting and 3DGS-based rendering. Here, \emph{markerless} refers to tracking
the demonstrator without worn markers or instrumented gloves; calibration
targets may still be used to establish camera geometry.

\section{Problem Statement}
\label{sec:problem}

The objective is to reuse a recorded human manipulation episode to generate
candidate robot observation--action sequences under several embodiment and
viewing configurations. Three connected requirements define the problem.
Perception must recover hand and object motion in a consistent metric frame.
Retargeting must adapt task-relevant motion to the robot's reach, joint limits,
end-effector geometry, and the geometric constraints represented in the model.
Rendering must produce observations consistent with the retargeted robot,
moving objects, and selected camera pose.

The intended expansion is described by $M$ robot embodiments, $K$ viewpoints,
and $P$ appearance configurations. Their Cartesian product contains up to
$M K P$ candidate episode configurations per source demonstration. It is not
a guarantee of $M K P$ feasible or independent demonstrations. Some robots may
not reach the task, some views may reveal unobserved surfaces, and some edits
may invalidate visual or geometric assumptions. The study therefore needs to
report both the requested combinations and the subset that passes explicit
checks.

Stage-level measurements alone are insufficient. A low hand-reconstruction
error, an admissible inverse-kinematics solution, and a plausible rendered
image can still disagree at their interfaces. The generated episode must also
preserve frame conventions, timestamps, object motion, and the correspondence
between rendered robot state and exported action. Geometric consistency does
not, by itself, establish force closure, successful manipulation, or usefulness
for policy learning.

\section{Research Gap}
\label{sec:gap}

The relevant literature includes integrated systems, not only isolated
components. RoboSplat augments a robot demonstration across object poses,
object types, camera views, scene appearance, lighting, and embodiments
\cite{robosplat}. Real2Render2Real extracts object motion from a human video
and generates robot data using reconstructed assets \cite{real2render2real}.
RoboPaint combines multiview RGB-D human capture, glove-based retargeting, and
hybrid rendering with a 3DGS background \cite{robopaint}. SplatSim generates
synthetic observations for simulated demonstrations, while GaussianFactory
synthesises demonstrations from a scene scan using grasp and trajectory
planning \cite{splatsim, gaussianfactory}.

Building on these systems, this thesis develops and assesses a particular
markerless multiview RGB-D route:
how reliably can reconstructed hand--object motion be reused across target
robots and rendered observation conditions, and where does this reuse fail?
The proposed contribution lies in the implementation and controlled
characterisation of that route, including the effects of capture geometry,
perception error, robot morphology, and rendering choices. Chapter~\ref{ch:lit}
compares the input requirements and representations of the closest systems.

\section{Research Questions}
\label{sec:rq}

The questions below define the intended evaluation of the pipeline. The
implementation and results in Chapter~\ref{ch:impl} provide the evidence for
answering them; questions without the required experiments remain open. A selected cohort of 28 HO-Cap manipulation sequences with eight
RGB-D views provides the source material for the perception assessment; it is
not the full dataset or its official held-out benchmark split.

\subsection*{RQ1: Hand Reconstruction Accuracy}

How accurately does landmark detection, multiview reconstruction, and temporal
regularisation recover human hand motion relative to the published HO-Cap
annotations?

\textit{Measures:} MPJPE, P95 joint error, PCK@20mm, and prediction coverage,
with directly observed and temporally filled predictions reported separately.

\textit{Reference:} HO-Cap's semi-automatic MANO-based annotations
\cite{hocap}. These are reference labels, not a competing baseline or an
independent motion-capture measurement. A parametric representation used by
the reference does not imply that ViKi fits the same model.

\subsection*{RQ2: Object Tracking Accuracy}

How accurately can rigid-object motion be recovered, and how does the result
depend on segmentation and initial pose information?

\textit{Measures:} Translation and rotation error, their aggregation level,
and tracking coverage. Orientation scoring must account for declared object
symmetries or report the limitations of the chosen pose convention.

\textit{Protocol:} Separate autonomous tracking from tracking with supplied
instance masks. Reference-assisted first-frame alignment measures subsequent
relative drift, not absolute localisation. Multiview agreement is a
self-consistency measure unless checked against an independent reference.

\subsection*{RQ3: Cross-Embodiment Retargeting Feasibility}

For which target robots can the recovered trajectories satisfy stated
kinematic and geometric constraints, and what tracking error remains?

\textit{Measures:} Accepted-episode and accepted-frame rates, TCP translation
and rotation error, joint and velocity-limit violations, self-collision
margin, and clearance from explicitly modelled scene geometry. Solver
convergence is reported separately from constraint satisfaction.

\textit{Candidate embodiments:} UR10, UR5, Franka Panda, Kinova Gen3, and
xArm7. The evaluation must state the actual robot models, base placements,
end-effectors, and trajectories tested. Kinematic acceptance does not certify
successful execution of the manipulation.

\subsection*{RQ4: Gripper State Transfer}

How consistently can observed thumb--index motion and interaction timing be
mapped to a parallel-jaw gripper command?

\textit{Measures:} Aperture tracking or command error against an explicitly
defined mapping, opening and closing timing relative to annotated events,
actuation-limit violations, and infeasible grasp-width cases. A mapped
command is not a measurement of achieved contact.

\textit{Scope:} Parallel-jaw transfer, initially using Robotiq 2F-85, is the
primary case. Robotiq 3F and Allegro Hand require richer hand correspondences
and a separate retargeting protocol; a scalar pinch signal does not determine
their complete joint configurations.

\subsection*{RQ5: Object-Aware Retargeting}

Does incorporating object motion improve geometric interaction consistency
relative to hand-only inverse kinematics?

\textit{Comparison:} Hand-only and object-relative retargeting on the same
source episodes and robot configurations; a contact-aware variant can be
included where a defined contact model is available.

\textit{Measures:} Relative hand/gripper--object pose error, unintended
penetration, and consistency of designated contact regions and timing.
These are geometric interaction measures, not physical grasp-success rates.
Replaying an imposed object trajectory cannot establish that the robot would
cause that motion.

\subsection*{RQ6: 3DGS Reconstruction and Novel-View Fidelity}

How do the number and placement of capture cameras affect reconstruction and
novel-view quality for the scene components represented by 3DGS?

\textit{Measures:} PSNR, SSIM, and valid-pixel depth RMSE on held-out views,
with separate evaluation of background and task-relevant object regions.

\textit{Protocol:} Fix train/test views and temporal splits before tuning.
The rendered reference must depict the same scene state. A static background
model must not be evaluated as if it reproduced the moving human and objects.
Timing perturbations, if evaluated, are a separate controlled ablation.

\subsection*{RQ7: Robot Compositing Consistency}

How consistently can a retargeted robot and moving objects be rendered within
the reconstructed scene across time and camera viewpoints?

\textit{Measures:} Depth-ordering and occlusion errors, boundary artifacts,
temporal stability, and agreement between rendered state and exported robot
state. Paired image metrics such as LPIPS require a corresponding reference
of the same target scene; a human demonstration frame is not a ground-truth
image of the replacement robot.

\textit{Comparison:} Alternative compositing methods on identical trajectories,
camera poses, and scene states. Human ratings, if collected, require a stated
sampling and rating protocol. Visual plausibility alone does not certify
correct interaction.

\subsection*{RQ8: Coverage and Yield of Dataset Expansion}

What coverage of embodiment, viewpoint, and appearance configurations can be
generated from each source demonstration while satisfying the defined
kinematic, geometric, and data-consistency checks?

\textit{Measures:} Requested and accepted episode counts, acceptance rate by
factor combination, viewpoint coverage, failure categories, and generation
time. Feature-space distances can describe visual variation under a fixed
encoder, but are not evidence of additional task skills or learning benefit.

\textit{Comparison:} Single-configuration export and controlled expansion of
the same source episodes. A teleoperation cost comparison requires measured
collection effort and an explicit accounting of calibration, scanning,
annotation, computation, and rejected outputs.

\medskip
\noindent The evaluation concerns reconstruction, retargeting, rendering, and
their integration. Policy training, task-success trials in simulation or on
hardware, and real-robot deployment are outside this study. Claims about
learning effectiveness or physically executable demonstrations require those
additional experiments.

\section{Intended Contributions}
\label{sec:contributions}

The planned contributions concern the implementation of the workflow and its
experimental characterisation:

\begin{enumerate}
  \item An offline pipeline design that separates multiview human-motion
        reconstruction, rigid-object tracking, robot retargeting, and visual
        synthesis through explicit intermediate representations.
  \item A workflow for reusing a source demonstration across robot,
        viewpoint, and appearance configurations, with rejected cases and
        configuration-dependent assumptions recorded rather than hidden in
        the nominal expansion count.
  \item An evaluation protocol for perception error, embodiment-specific
        feasibility, novel-view fidelity, and observation--action consistency,
        including controlled comparisons where reference data are available.
  \item A characterisation of the resulting dataset coverage and generation
        effort within the tested tasks and embodiments.
\end{enumerate}

The intended outcome is a dataset-generation workflow with measured operating
limits, acceptance criteria, and documented failure cases.

\section{Thesis Structure}
\label{sec:structure}

Chapter~\ref{ch:lit} reviews human-motion representations, retargeting,
3DGS-based data generation, task representations, and capture geometry.
Chapter~\ref{ch:method} describes the pipeline formulations and assumptions.
Chapter~\ref{ch:impl} reports the implementation and available evaluation.
Chapter~\ref{ch:disc} examines measurement validity, failure cases, and the
limits of generalisation. Chapter~\ref{ch:concl} summarises the supported
findings and identifies the questions that require further experiments.
